% Section 12.3: The exterior derivative | Part: 2 | Depends on: sec:tangent, sec:ext | Status: done | Original lines: 2056-2145
%!TEX root = ../../main.tex
\subsection{The exterior derivative}\label{sec:dform}
% REWRITTEN: the repeated introduction of T_pM, \partial_i and dx^i is replaced by a pointer to Section sec:tangentvectors.
Let $M$ be a smooth $n$-manifold (for concreteness, an open subset of $\R^n$ with coordinates $x^1,\dots,x^n$), with the coordinate bases
$\partial_i$ of $T_pM$ and $\dd x^i$ of $T_p^*M$ of Section~\ref{sec:tangentvectors}.

\begin{definition}[Differential form]
A \emph{differential $k$-form} is a smooth assignment $p\mapsto\omega_p\in\bigwedge^kT_p^*M\cong\Alt^k(T_pM;\R)$, i.e.\ at each point an alternating $k$-linear function on tangent vectors (Remark~\ref{rem:twomodels}). In coordinates
\[
 \omega=\sum_{i_1<\cdots<i_k}\omega_{i_1\cdots i_k}(x)\,\dd x^{i_1}\wedge\cdots\wedge\dd x^{i_k}=\sum_I\omega_I\,\dd x^I,
\]
with smooth coefficients. $\Omega^k(M)$ denotes the space of $k$-forms; $0$-forms are functions. The wedge product is the pointwise exterior product.
\end{definition}

\begin{definition}[Exterior derivative]\label{def:d}
The \emph{exterior derivative} is the unique $\R$-linear operator $\dd:\Omega^k(M)\to\Omega^{k+1}(M)$ such that
\begin{enumerate}
  \item for functions, $\dd f=\sum_i\partial_if\,\dd x^i$ (i.e.\ $\dd f(X)=Xf$);
  \item (graded Leibniz rule) $\dd(\alpha\wedge\beta)=\dd\alpha\wedge\beta+(-1)^k\alpha\wedge\dd\beta$ for $\alpha\in\Omega^k$;
  \item $\dd^2=0$.
\end{enumerate}
In coordinates:
\begin{equation}\label{eq:dcoord}
 \dd\Bigl(\sum_I\omega_I\,\dd x^I\Bigr)=\sum_I\sum_{j=1}^n\partial_j\omega_I\;\dd x^j\wedge\dd x^I .
\end{equation}
\end{definition}

% REWRITTEN (F4): the Lie bracket used in (3) is defined inline (it is studied in Section sec:lie); (4) is marked as a
% preview of Theorem thm:stokes; pull-back written \phi\pb instead of \phi^{*}.
% AUTHOR-NOTE: (2), (3) and (5) are stated without proof here. Glossary Part C proves (2) (id:pbd) and (3) for 1-forms
% (id:dinvariant); (3) for k>1 and (5) (Poincare lemma) are left to the references (docs/AUTHOR_NOTES.md #7).
\begin{theorem}[Fundamental properties]\label{thm:dprops}
\begin{enumerate}
  \item $\dd^2=0$.
  \item $\dd$ commutes with pullbacks: $\phi\pb(\dd\omega)=\dd(\phi\pb\omega)$ for every smooth map $\phi$. In particular formula \eqref{eq:dcoord} is independent of coordinates.
  \item (\emph{Invariant formula}) For vector fields $X_0,\dots,X_k$, where $[X,Y]$ denotes the \emph{Lie bracket}, the vector field
  $[X,Y]f=X(Yf)-Y(Xf)$ (Section~\ref{sec:lie}):
  \begin{multline*}
  \dd\omega(X_0,\dots,X_k)=\sum_i(-1)^iX_i\bigl(\omega(\dots,\widehat{X_i},\dots)\bigr)\\
  +\sum_{i<j}(-1)^{i+j}\omega([X_i,X_j],\dots,\widehat{X_i},\dots,\widehat{X_j},\dots).
  \end{multline*}
  \item (\emph{Stokes' theorem}, preview of Theorem~\ref{thm:stokes}, where orientation, boundary and the integral of a form are defined) If $M$ is an oriented compact $n$-manifold with boundary and $\omega\in\Omega^{n-1}(M)$, then $\displaystyle\int_M\dd\omega=\int_{\partial M}\omega$.
  \item (\emph{Poincar\'e lemma}) On a contractible (e.g.\ star-shaped) domain, every closed form ($\dd\omega=0$) of degree $\ge1$ is exact ($\omega=\dd\eta$).
\end{enumerate}
\end{theorem}
\begin{proof}[Proof of (1)]
By \eqref{eq:dcoord} applied twice,
$\dd^2\omega=\sum_I\sum_{j,k}\partial_k\partial_j\omega_I\,\dd x^k\wedge\dd x^j\wedge\dd x^I.$
The coefficient $\partial_k\partial_j\omega_I$ is symmetric in $(j,k)$ (equality of mixed partials), whereas $\dd x^k\wedge\dd x^j$ is antisymmetric; so the sum vanishes.
\end{proof}
% REWRITTEN: glossary pointer
Proofs of (2), and of (3) for $1$-forms, are Identities~\ref{id:pbd} and~\ref{id:dinvariant} in the Glossary; for (3) in general and for (5) see \cite{Lee,BottTu}.

\begin{example}[$\R^3$: gradient, curl, divergence]
\begin{itemize}
  \item \emph{gradient} ($0$-forms $\to$ $1$-forms):
  $\dd f=f_x\,\dd x+f_y\,\dd y+f_z\,\dd z$;
  \item \emph{curl} ($1$-forms $\to$ $2$-forms):
  \begin{align*}
   \dd(P\,\dd x+Q\,\dd y+R\,\dd z)&=(R_y-Q_z)\,\dd y\wedge\dd z+(P_z-R_x)\,\dd z\wedge\dd x\\
   &\quad+(Q_x-P_y)\,\dd x\wedge\dd y;
  \end{align*}
  \item \emph{divergence} ($2$-forms $\to$ $3$-forms):
  $\dd(A\,\dd y\wedge\dd z+B\,\dd z\wedge\dd x+C\,\dd x\wedge\dd y)=(A_x+B_y+C_z)\,\dd x\wedge\dd y\wedge\dd z$.
\end{itemize}

Thus $\dd^2=0$ contains the identities $\vd\times\vd f=0$ and $\vd\cdot(\vd\times\mathbf v)=0$, and Stokes' theorem contains the gradient theorem, the Kelvin--Stokes theorem and the divergence theorem.
\end{example}

% REORDERED (F5): the step-by-step verification of curl and divergence was in Section sec:classical (original L2798-2812);
% it uses only d, wedge and eq:iotavol, and is needed by Synthesis II. The integral theorems (iii), (iv) stay in sec:classical.
\subhead{Curl and divergence as $\dd$, step by step}
In $\R^3$ with the Euclidean metric and volume form $\vol=\dd x\wedge\dd y\wedge\dd z$, lowering an index changes nothing numerically: a vector field $F$ gives $F^\flat=F_x\dd x+F_y\dd y+F_z\dd z$. Recall \eqref{eq:iotavol}: $\iota_G\vol=G_x\,\dd y\wedge\dd z+G_y\,\dd z\wedge\dd x+G_z\,\dd x\wedge\dd y$.

\emph{(i) Curl.} Compute $\dd F^\flat$ term by term, using $\dd x\wedge\dd x=0$:
\begin{align*}
\dd(F_x\dd x)&=\partial_yF_x\,\dd y\wedge\dd x+\partial_zF_x\,\dd z\wedge\dd x=-\partial_yF_x\,\dd x\wedge\dd y+\partial_zF_x\,\dd z\wedge\dd x,\\
\dd(F_y\dd y)&=\partial_xF_y\,\dd x\wedge\dd y+\partial_zF_y\,\dd z\wedge\dd y=\partial_xF_y\,\dd x\wedge\dd y-\partial_zF_y\,\dd y\wedge\dd z,\\
\dd(F_z\dd z)&=\partial_xF_z\,\dd x\wedge\dd z+\partial_yF_z\,\dd y\wedge\dd z=-\partial_xF_z\,\dd z\wedge\dd x+\partial_yF_z\,\dd y\wedge\dd z.
\end{align*}
Adding and collecting the three basis $2$-forms:
\[
\dd F^\flat=(\partial_yF_z-\partial_zF_y)\,\dd y\wedge\dd z+(\partial_zF_x-\partial_xF_z)\,\dd z\wedge\dd x+(\partial_xF_y-\partial_yF_x)\,\dd x\wedge\dd y=\iota_{\operatorname{curl}F}\vol,
\]
since the three coefficients are exactly the components of $\operatorname{curl}F$. Because $\hodge\hodge=1$ in $\R^3$ and $\hodge(\dd y\wedge\dd z)=\dd x$, etc.\ (Section~\ref{sec:duality}), this is the same as $\operatorname{curl}F=(\hodge\,\dd F^\flat)^\sharp$.

\emph{(ii) Divergence.} $\iota_F\vol=F_x\dd y\wedge\dd z+F_y\dd z\wedge\dd x+F_z\dd x\wedge\dd y$. Differentiating, only the derivative with respect to the variable missing from each $2$-form survives:
$\dd(F_x\dd y\wedge\dd z)=\partial_xF_x\,\dd x\wedge\dd y\wedge\dd z$, $\dd(F_y\dd z\wedge\dd x)=\partial_yF_y\,\dd y\wedge\dd z\wedge\dd x$, $\dd(F_z\dd x\wedge\dd y)=\partial_zF_z\,\dd z\wedge\dd x\wedge\dd y$,
and the three $3$-forms equal $\dd x\wedge\dd y\wedge\dd z$ (cyclic permutations of three $1$-forms are even). So $\dd(\iota_F\vol)=(\operatorname{div}F)\,\vol$.

\begin{intuition}
\begin{itemize}
  \item A $1$-form is a field of ``measuring devices'' for line segments: integrate it along a path to get a number (work, circulation). A $2$-form measures oriented \emph{areas} (flux), a $3$-form measures volumes. This is exactly the dual (``measuring'') side of the arrows/areas picture of \S\ref{sec:ext}.
  \item $\dd\omega$ is the \emph{infinitesimal circulation}: $\dd\omega(X,Y)\,\epsilon^2$ is the integral of $\omega$ around the little parallelogram spanned by $\epsilon X,\epsilon Y$ (the closing term of the fundamental theorem of calculus). Stokes' theorem is the statement that if you sum these small circulations over a region, interior edges cancel in pairs and only the boundary survives.
  \item $\dd^2=0$ is the algebraic shadow of ``\emph{the boundary of a boundary is empty}'': the faces of an infinitesimal cube have boundary edges that cancel in pairs.
  \item Grad, curl and div are \emph{one} operator $\dd$ acting in different degrees; the three vector-calculus theorems are \emph{one} theorem. The cross product and the ``vector'' character of curl are artefacts of $3$D, where $2$-forms can be disguised as vectors by duality.
\end{itemize}
\end{intuition}


\begin{example}[Everything on one function]\label{ex:onefunction}
Let $f(x,y)=x^2y$ on $\R^2$. Then
\begin{align*}
\dd f&=2xy\,\dd x+x^2\,\dd y &&\text{(second relation of \eqref{eq:trio})}\\
\partial_xf&=\iota_{\partial_x}\dd f=\dd f(\partial_x)=2xy &&\text{(third relation of \eqref{eq:trio})}\\
\dd(\dd f)&=\dd(2xy)\wedge\dd x+\dd(x^2)\wedge\dd y &&\text{(Leibniz, $\dd\dd x=\dd\dd y=0$)}\\
&=(2y\,\dd x+2x\,\dd y)\wedge\dd x+2x\,\dd x\wedge\dd y\\
&=2x\,\dd y\wedge\dd x+2x\,\dd x\wedge\dd y\\
&=-2x\,\dd x\wedge\dd y+2x\,\dd x\wedge\dd y=0 &&\text{(antisymmetry).}
\end{align*}
The last line is $\dd\dd=0$: it holds because $\partial_x\partial_yf=\partial_y\partial_xf$ is symmetric while $\dd x\wedge\dd y$ is antisymmetric (Theorem~\ref{thm:dprops}(1)).
% REWRITTEN: cited thm:dprops (backward) instead of sec:expl-closed (forward)
\end{example}

\begin{keyidea}[$\dd$ needs no metric and no connection]
The exterior derivative is built from partial derivatives and the wedge product only. Its coordinate independence (Theorem~\ref{thm:dprops}(2)) comes from the symmetry of mixed partials against the antisymmetry of $\wedge$.
In Section~\ref{sec:synth3} we shall see why a covariant derivative, whose second derivatives do \emph{not} commute (curvature), cannot be used to build such an operator unless the space is flat.
\end{keyidea}
